%==========================================================
%  CHAPTER 4: Labour Markets — Employment, Wages, and
%              Demographic Change
%  ECON 204 — Contemporary Economic Issues: A Canadian Perspective
%==========================================================

\chapter{Labour Markets: Employment, Wages, and Demographic Change}
\label{ch:labour}

%----------------------------------------------------------
%  OPENING VIGNETTE
%----------------------------------------------------------

\noindent In March 2023, two Canadians were navigating very different
labour markets --- in the same city.

Jerome, 42, had worked as a machinist at a Windsor, Ontario auto parts
plant for eighteen years. His unionized position paid \$34.20 per hour,
came with a defined-benefit pension, dental and vision coverage, and
four weeks of paid vacation. He had renegotiated his collective agreement
the previous fall and received a 6.5\% wage increase --- his largest in
a decade. His plant had just won a \$200 million retooling contract for
electric vehicle battery components. His employer was begging for more
skilled machinists and couldn't find them.

Kayla, 27, had a business degree from the University of Windsor and
worked as a ``gig coordinator'' for a logistics platform, classifying
her own work schedule and delivery routes through an app. Her pay was
\$19.50 per effective hour after accounting for vehicle costs, which
she covered herself. She had no benefits, no pension, no vacation pay,
and no job security beyond the next available delivery slot. On weeks when
demand was low, her earnings fell below minimum wage. On weeks when it was
high --- Christmas, Valentine's Day --- she worked 60 hours.

Jerome and Kayla live in the same city, are subject to the same provincial
employment standards, and are affected by the same monetary and fiscal
policies. Yet their labour market experiences are separated by a chasm:
one navigated by institutional arrangements built over decades of
collective bargaining; the other by algorithms, platform economics, and
legal classifications that exclude workers from the protections most
Canadians take for granted.

Understanding labour markets --- how wages are determined, why employment
fluctuates, how technological change and demographic aging are reshaping
the demand for and supply of work, and what policies can improve outcomes
for workers like Kayla without undermining conditions for workers like
Jerome --- is the subject of this chapter.

%----------------------------------------------------------
%  LEARNING OBJECTIVES
%----------------------------------------------------------

\begin{learningobjectives}
  \item Apply the supply-and-demand model to labour markets, distinguish
        the competitive from the monopsonistic wage-setting case, and
        use both to predict the employment effects of a minimum wage.
  \item Calculate the three key labour force statistics --- unemployment
        rate, employment rate, and labour force participation rate ---
        from raw population and employment data, and interpret each
        measure correctly.
  \item Distinguish among frictional, structural, and cyclical
        unemployment, define the natural rate of unemployment, and
        explain why the unemployment rate is an imperfect welfare measure.
  \item Interpret the Beveridge Curve and explain what an outward shift
        of the curve means for labour market matching efficiency.
  \item Analyse the wage-productivity gap in Canada, identify its causes,
        and evaluate evidence on the link between unionization decline
        and stagnant real wages.
  \item Compare standard and non-standard employment on wages, benefits,
        and security, and evaluate the policy debate over the employment
        status of platform workers.
  \item Assess the demographic pressures on Canada's labour market ---
        aging population, declining youth cohorts, and immigration ---
        and connect them to labour shortages, the dependency ratio, and
        long-run productivity.
\end{learningobjectives}

%==========================================================
\section{Labour Market Fundamentals}
\label{sec:labour_fundamentals}
%==========================================================

The \term{labour market} is the market in which workers sell their time
and skills to employers. Like any market, it has a supply side (workers
deciding how many hours to offer at different wage rates), a demand side
(employers deciding how many workers to hire at different wage rates), and
an equilibrium (the wage and employment level that clears the market).

\begin{defbox}{Labour Demand and Labour Supply}
\textbf{Labour demand} is derived from the demand for output.
A firm's profit-maximising condition for labour hiring is:
\[
  W = P \times \text{MP}_L
\]
where $W$ is the market wage, $P$ is the output price, and $\text{MP}_L$
is the marginal product of labour. This condition equates the wage to
the \term{value of the marginal product of labour (VMP$_L$)}, which is
the revenue generated by the last unit of labour hired. Labour demand
curves slope downward because, due to diminishing returns, each
additional worker adds less output (and hence less revenue) than the
previous one.

\textbf{Labour supply} reflects the trade-off between work (earning
income) and leisure (consuming non-work time). The supply curve
typically slopes upward: higher wages induce more labour supply,
as the opportunity cost of leisure rises. However, at very high wages,
the \term{income effect} (workers can afford more leisure) can dominate
the \term{substitution effect} (higher wage makes leisure more costly),
producing a backward-bending supply curve.
\end{defbox}

%----------------------------------------------------------
\subsection{The Competitive Labour Market}
%----------------------------------------------------------

In a perfectly competitive labour market, many workers and many firms
interact, no individual can influence the market wage, and the equilibrium
wage $W^*$ clears the market at employment $L^*$. This is the benchmark
model. Figure~\ref{fig:labour_markets} illustrates both the competitive
case and its main departure: monopsony.

\begin{figure}[H]
\centering
\begin{tikzpicture}[
  axisline/.style={-Stealth, thick, color=gray!70},
  lbl/.style={font=\small}
]

%---- PANEL A: COMPETITIVE MARKET ----
\begin{scope}[xshift=0cm]
\node[font=\small\bfseries, align=center] at (3.2,5.8)
  {Panel A: Competitive Market};
\draw[axisline] (0,0) -- (0,5.5) node[above, lbl] {$W$};
\draw[axisline] (0,0) -- (6.0,0) node[right, lbl] {$L$};
% Supply
\draw[casered, thick] (0.4,0.4) -- (5.4,4.9)
  node[right, font=\footnotesize, color=casered] {$S_L$};
% Demand
\draw[policyblue, thick] (0.4,4.9) -- (5.6,0.3)
  node[right, font=\footnotesize, color=policyblue] {$D_L = \text{VMP}_L$};
% Equilibrium
\fill[unbcgreen] (2.88,2.68) circle (3pt);
\draw[dashed, gray!50, thin] (0,2.68) -- (2.88,2.68)
  node[left, pos=0, font=\footnotesize] {$W^*$};
\draw[dashed, gray!50, thin] (2.88,0) -- (2.88,2.68)
  node[below, font=\footnotesize] {$L^*$};
% Minimum wage
\draw[dataorange, thick, dashed] (0,3.8) -- (5.8,3.8)
  node[right, font=\footnotesize, color=dataorange] {$\bar{W}_{\min}$};
\fill[dataorange] (1.62,3.8) circle (2.5pt);
\fill[dataorange] (4.35,3.8) circle (2.5pt);
\draw[<->, dataorange!80, thin] (1.62,4.1) -- (4.35,4.1)
  node[midway, above, font=\footnotesize, color=dataorange] {unemployment};
\node[left, font=\footnotesize, dataorange] at (0,3.8) {$\bar{W}$};
\end{scope}

%---- PANEL B: MONOPSONY ----
\begin{scope}[xshift=7.5cm]
\node[font=\small\bfseries, align=center] at (3.2,5.8)
  {Panel B: Monopsony};
\draw[axisline] (0,0) -- (0,5.5) node[above, lbl] {$W$};
\draw[axisline] (0,0) -- (6.0,0) node[right, lbl] {$L$};
% Supply (ACL)
\draw[casered, thick] (0.4,0.4) -- (5.4,4.9)
  node[right, font=\footnotesize, color=casered] {$S_L = \text{ACL}$};
% MLC (steeper than supply)
\draw[casered!60, thick, dashed] (0.4,0.8) -- (3.8,5.2)
  node[right, font=\footnotesize, color=casered!70] {MLC};
% VMP_L (demand)
\draw[policyblue, thick] (0.4,4.9) -- (5.6,0.3)
  node[right, font=\footnotesize, color=policyblue] {$D_L = \text{VMP}_L$};
% Monopsony equilibrium
\fill[casered!80] (2.15,2.15) circle (3pt);
% Monopsony wage (read off supply curve)
\draw[dashed, gray!50, thin] (0,1.55) -- (2.15,1.55)
  node[left, pos=0, font=\footnotesize] {$W_m$};
\draw[dashed, gray!50, thin] (2.15,0) -- (2.15,2.15)
  node[below, font=\footnotesize] {$L_m$};
% Competitive outcome for comparison
\fill[unbcgreen] (2.88,2.68) circle (2.5pt);
\draw[dashed, gray!20, thin] (0,2.68) -- (2.88,2.68)
  node[left, pos=0, font=\footnotesize, color=gray!60] {$W^*$};
\draw[dashed, gray!20, thin] (2.88,0) -- (2.88,2.68)
  node[below, font=\footnotesize, color=gray!60] {$L^*$};
% Minimum wage eliminating monopsony distortion
\draw[dataorange, thick, dashed] (0,2.68) -- (5.8,2.68)
  node[right, font=\footnotesize, color=dataorange] {$\bar{W}_{\min} = W^*$};
\node[font=\footnotesize, align=center, text=dataorange] at (4.2,1.3)
  {Min wage restores\\competitive outcome};
\end{scope}
\end{tikzpicture}
\caption{Labour Market Models: Competitive vs.\ Monopsony.
\textbf{Panel A} shows the competitive labour market. The minimum wage
$\bar{W}$ set above the competitive equilibrium $W^*$ creates unemployment
(excess labour supply, shown by the orange gap) --- the standard result.
\textbf{Panel B} shows a monopsony (single buyer of labour). The monopsonist
sets employment where MLC $=$ VMP$_L$, at $L_m < L^*$, and pays wage
$W_m < W^*$. Critically, a minimum wage set at $W^*$ eliminates the
monopsony distortion and \textit{increases} both employment and wages ---
the theoretical basis for why Card and Krueger (1994) found no employment
loss from minimum wage increases in some US fast-food markets. Canadian
labour markets, particularly in rural communities and in care sectors,
often have monopsonistic features.}
\label{fig:labour_markets}
\end{figure}

%----------------------------------------------------------
\subsection{Monopsony in Canadian Labour Markets}
%----------------------------------------------------------

The competitive model assumes workers can freely move between employers.
In practice, many Canadian workers face a \term{monopsonist} or
\term{oligopsonist} labour market --- a market where one or a few employers
dominate hiring. Classic examples include:

\begin{itemize}
  \item \textbf{Single-employer towns:} A sawmill in a remote Northern BC
        community, a mine in the Northwest Territories, or a fish processing
        plant in rural Newfoundland may be the only significant employer.
        Workers who do not want to relocate have little bargaining power.
  \item \textbf{Sector-specific skills:} A highly trained nurse or
        schoolteacher in a small city faces limited employer options
        even if the labour market is otherwise competitive.
  \item \textbf{Platform labour markets:} Ride-share and delivery platforms
        effectively set wages unilaterally for workers classified as
        contractors --- exhibiting monopsonistic wage-setting in practice
        even when no legal monopoly exists.
\end{itemize}

The monopsony model has an important policy implication: when labour markets
are monopsonistic, a minimum wage set up to the competitive wage level can
\textit{increase} employment and wages simultaneously. This is why the
empirical evidence on minimum wages is more nuanced than the simple
competitive model predicts. We return to this in Section~\ref{sec:min_wage}.

%==========================================================
\section{Labour Force Accounting}
\label{sec:labour_accounting}
%==========================================================

The Statistics Canada \term{Labour Force Survey (LFS)}, published monthly,
is Canada's primary source of employment and unemployment data. The LFS
divides the working-age population (WAP, persons aged 15 and over) into
three groups:

\begin{itemize}
  \item \textbf{Employed} ($E$): persons who worked at least one hour
        for pay or profit in the reference week, or who were absent from
        a job they held (due to illness, vacation, or labour dispute).
  \item \textbf{Unemployed} ($U$): persons who were without work, were
        available for work, and were actively seeking employment in the
        four weeks prior to the reference week.
  \item \textbf{Not in the labour force} ($N$): persons who are neither
        employed nor unemployed --- including retirees, students, caregivers,
        and discouraged workers who have stopped searching for jobs.
\end{itemize}

The labour force is $LF = E + U$. From these groups, three key statistics
are derived:

\begin{equation}
  \text{Unemployment Rate} = \frac{U}{LF} \times 100 = \frac{U}{E + U} \times 100
  \label{eq:unemployment_rate}
\end{equation}

\begin{equation}
  \text{Labour Force Participation Rate (LFPR)} = \frac{LF}{\text{WAP}} \times 100
  \label{eq:lfpr}
\end{equation}

\begin{equation}
  \text{Employment Rate} = \frac{E}{\text{WAP}} \times 100
  \label{eq:employment_rate}
\end{equation}

\noindent\textbf{Worked Example 4.1 --- Labour Force Accounting, Canada, Mid-2024}

\begin{center}
\renewcommand{\arraystretch}{1.35}
\begin{tabular}{lr}
\toprule
\textbf{Category} & \textbf{Value} \\
\midrule
Working-age population (WAP, 15+)             & 34{,}150{,}000 \\
Employed ($E$)                                & 20{,}615{,}000 \\
Unemployed ($U$)                              &  1{,}371{,}000 \\
Not in labour force ($N$)                     & 12{,}164{,}000 \\
\midrule
Labour force ($LF = E + U$)                   & 21{,}986{,}000 \\
\midrule
Unemployment rate ($U/LF$)                    & 6.2\% \\
Labour force participation rate ($LF$/WAP)    & 64.4\% \\
Employment rate ($E$/WAP)                     & 60.4\% \\
\bottomrule
\end{tabular}
\end{center}

\noindent\textit{Interpretation:}
\begin{itemize}
  \item A 6.2\% unemployment rate means approximately 1.37 million
        Canadians who want to work and are actively searching cannot
        find a job.
  \item A 64.4\% LFPR means that 35.6\% of working-age Canadians are
        neither employed nor looking for work.
  \item An employment rate of 60.4\% means only 60 of every 100 Canadians
        aged 15+ hold a paying job.
\end{itemize}

\begin{keytakeaway}
The unemployment rate is the most cited labour market statistic, but it
is an imperfect welfare measure. It excludes: \textbf{discouraged workers}
(who want work but stopped searching); \textbf{involuntary part-timers}
(who want full-time work but can only find part-time); and \textbf{workers
in jobs below their skill level}. Statistics Canada's R8 measure of
underemployment, which includes these groups, consistently runs 3--5
percentage points above the headline unemployment rate.
\end{keytakeaway}

%----------------------------------------------------------
\subsection{Canada's Labour Market Record, 2010--2026}
%----------------------------------------------------------

\begin{figure}[H]
\centering
\begin{tikzpicture}
\begin{axis}[
  width=14cm, height=7.5cm,
  xlabel={Year},
  ylabel={Rate (\%)},
  xlabel style={font=\small},
  ylabel style={font=\small},
  xmin=2009.5, xmax=2026.5,
  ymin=0, ymax=75,
  xtick={2010,2012,2014,2016,2018,2020,2022,2024,2026},
  xticklabel style={font=\small},
  ytick={0,10,20,30,40,50,60,70},
  yticklabel style={font=\small},
  ymajorgrids=true,
  grid style={dashed, gray!20},
  legend pos=south east,
  legend style={font=\small},
  clip=false
]
% LFPR
\addplot[policyblue, thick, mark=none] coordinates {
  (2010,67.1)(2011,66.8)(2012,66.6)(2013,66.5)(2014,65.8)
  (2015,65.9)(2016,65.7)(2017,65.8)(2018,65.5)(2019,65.6)
  (2020,62.4)(2021,64.7)(2022,64.8)(2023,65.1)(2024,64.4)(2025,64.1)(2026,64.3)
};
\addlegendentry{LFPR (\%, left)}
% Employment rate
\addplot[unbcgreen, thick, dashed, mark=none] coordinates {
  (2010,61.6)(2011,61.7)(2012,61.8)(2013,62.0)(2014,61.4)
  (2015,61.4)(2016,61.0)(2017,61.6)(2018,62.0)(2019,61.8)
  (2020,56.5)(2021,59.9)(2022,61.5)(2023,61.3)(2024,60.4)(2025,59.9)(2026,60.1)
};
\addlegendentry{Employment rate (\%)}
% Unemployment rate (scaled up ×5 for visibility on same axis)
\addplot[casered, thick, dotted, mark=none] coordinates {
  (2010,8.1)(2011,7.4)(2012,7.2)(2013,7.1)(2014,6.9)
  (2015,6.9)(2016,7.0)(2017,6.3)(2018,5.8)(2019,5.7)
  (2020,9.5)(2021,7.5)(2022,5.3)(2023,5.4)(2024,6.3)(2025,6.7)(2026,6.5)
};
\addlegendentry{Unemployment rate (\%)}
% COVID annotation
\node[font=\footnotesize, align=center, color=gray] at (axis cs:2020,50)
  {COVID-19\\lockdowns};
\draw[->, gray!60, thin] (axis cs:2020,52) -- (axis cs:2020,57);
\node[font=\footnotesize, color=casered] at (axis cs:2020,11)
  {u=9.5\%};
\node[font=\footnotesize, color=casered] at (axis cs:2022.5,4.5)
  {u=5.3\%\\(50-yr low)};
\end{axis}
\end{tikzpicture}
\caption{Canadian Labour Force Participation Rate, Employment Rate, and
Unemployment Rate, 2010--2026. The three series together tell a richer
story than any one alone. The COVID-19 collapse of April 2020 is visible
in all three: the unemployment rate spiked to 9.5\%, the employment rate
fell to 56.5\%, and the LFPR dropped to 62.4\% as many workers stopped
searching. The 2022 labour market tightness is equally clear: unemployment
at 4.9--5.3\% (its lowest since the 1970s) while LFPR had recovered to
near pre-pandemic levels. The subsequent 2024--2025 softening reflects
the economic slowdown induced by the Bank of Canada's rate-hike cycle
(Chapter~\ref{ch:inflation}).}
\label{fig:lf_statistics}
\source{Statistics Canada, Table 14-10-0287-01, Labour Force Survey,
seasonally adjusted. 2025--2026 estimated.}
\end{figure}

%==========================================================
\section{Types of Unemployment}
\label{sec:unemployment_types}
%==========================================================

Not all unemployment is the same. Economists distinguish three conceptually
distinct types, each with different causes and policy implications.

\begin{defbox}{Three Types of Unemployment}
\textbf{Frictional unemployment} arises from the normal time it takes
to match workers with jobs. Even in a healthy economy, workers quit,
are dismissed, or leave school and must search for the right job.
Firms take time to evaluate candidates. This process generates temporary
unemployment that is voluntary, transitional, and inevitable in any
dynamic market economy. Estimated at 1--2\% of the labour force.

\textbf{Structural unemployment} arises when workers' skills are
mismatched with available jobs --- either because of technological
change (automation displacing specific occupations), geographic
mismatch (jobs in growing cities, workers in declining regions), or
industry restructuring (oil sands declining while green tech expands).
Structural unemployment tends to be longer-duration and is not cured
by aggregate demand stimulus. Policy responses include retraining,
relocation assistance, and apprenticeship programs.

\textbf{Cyclical unemployment} arises from insufficient aggregate
demand during economic downturns. When GDP falls below potential,
firms reduce hiring and may lay off workers; the resulting unemployment
is correlated with the business cycle. Cyclical unemployment is the
target of countercyclical monetary and fiscal policy.

The \textbf{natural rate of unemployment} $u_n$ (also called the
NAIRU) is the unemployment rate that prevails when the economy is
at potential output, consisting of frictional and structural
unemployment only. For Canada, the Bank of Canada estimates
$u_n \approx 5.5$--$6.0\%$ in the mid-2020s.
\end{defbox}

The distinction between structural and cyclical unemployment matters
enormously for policy. When the Bank of Canada raised rates in 2022--2023,
it was targeting the \textit{cyclical} component of inflation (excess
demand). The resulting rise in unemployment from 5.3\% to 6.5--6.8\%
by 2024--25 reflects both cyclical slack (deliberate demand reduction)
and structural factors (immigration increasing labour supply faster than
job creation). A policymaker who mistakes structural unemployment for
cyclical --- and responds with rate cuts --- risks reflating the economy
before the structural adjustment is complete.

%==========================================================
\section{The Beveridge Curve: Labour Market Matching}
\label{sec:beveridge_curve}
%==========================================================

The \term{Beveridge Curve} plots the job vacancy rate (share of positions
unfilled) against the unemployment rate over time. It describes the
efficiency of the process by which workers and jobs are matched.

A well-functioning labour market traces a stable, downward-sloping curve:
when unemployment is high (recession), vacancies are low (few firms
hiring); when unemployment is low (tight market), vacancies are high
(firms competing intensely for scarce workers). The position of the curve
reflects \term{matching efficiency} --- how well the process of connecting
workers to jobs works.

\begin{figure}[H]
\centering
\begin{tikzpicture}
\begin{axis}[
  width=10cm, height=9cm,
  xlabel={Unemployment Rate (\%)},
  ylabel={Job Vacancy Rate (\%)},
  xlabel style={font=\small},
  ylabel style={font=\small},
  xmin=4.5, xmax=10.5,
  ymin=0.5, ymax=5.5,
  xtick={5,6,7,8,9,10},
  ytick={1,2,3,4,5},
  xticklabel style={font=\small},
  yticklabel style={font=\small},
  ymajorgrids=true, xmajorgrids=true,
  grid style={dashed, gray!20},
  clip=false
]
% Pre-COVID Beveridge curve (reference)
\addplot[gray!40, dashed, thick, domain=4.5:10.5]
  {1.8 + 0.8*exp(-0.45*(x-5.0))}
  node[above right, pos=0.05, font=\footnotesize, color=gray!60]
  {Pre-COVID curve};
% Post-COVID Beveridge curve (outward shift)
\addplot[unbcgreen!50, dashed, thick, domain=4.5:10.5]
  {2.5 + 1.2*exp(-0.45*(x-5.0))}
  node[above right, pos=0.05, font=\footnotesize, color=unbcgreen!70]
  {Post-COVID curve};
% Data points
\addplot[only marks, mark=*, mark size=3pt, policyblue]
  coordinates {
    (6.9,1.5)(7.0,1.6)(6.3,1.9)(5.8,2.1)(5.7,2.2)
  };
\addplot[only marks, mark=square*, mark size=3pt, casered]
  coordinates {
    (9.5,1.4)(7.5,3.2)(5.3,4.5)(5.4,3.4)(6.3,2.8)(6.7,2.5)
  };
% Year labels
\node[font=\footnotesize, policyblue, above right] at (axis cs:6.9,1.5) {15};
\node[font=\footnotesize, policyblue, above right] at (axis cs:7.0,1.6) {16};
\node[font=\footnotesize, policyblue, above right] at (axis cs:6.3,1.9) {17};
\node[font=\footnotesize, policyblue, above right] at (axis cs:5.8,2.1) {18};
\node[font=\footnotesize, policyblue, above right] at (axis cs:5.7,2.2) {19};
\node[font=\footnotesize, casered, right] at (axis cs:9.5,1.4) {20};
\node[font=\footnotesize, casered, above right] at (axis cs:7.5,3.2) {21};
\node[font=\footnotesize, casered, above] at (axis cs:5.3,4.5) {22};
\node[font=\footnotesize, casered, right] at (axis cs:5.4,3.4) {23};
\node[font=\footnotesize, casered, below right] at (axis cs:6.3,2.8) {24};
\node[font=\footnotesize, casered, below right] at (axis cs:6.7,2.5) {25};
% Arrow showing outward shift
\draw[->, unbcgreen!70, thick] (axis cs:6.5,2.3) -- (axis cs:6.3,3.1);
\node[font=\footnotesize, unbcgreen!80, align=center] at (axis cs:7.0,3.7)
  {Outward\\shift:\\matching\\inefficiency};
\end{axis}
\end{tikzpicture}
\caption{Beveridge Curve: Canada, 2015--2025. Blue circles (labelled
15--19) show the pre-pandemic period tracing a stable downward-sloping
relationship near the grey reference curve. Red squares (labelled 20--25)
show the post-COVID period. The 2021--2022 points sit far to the
upper-left of the pre-COVID curve: very high vacancies (4.5\% in 2022)
combined with still-elevated unemployment (5.3\%), suggesting that many
unfilled positions could not be matched with available workers. This
\textit{outward shift} of the Beveridge curve indicates reduced matching
efficiency --- caused by skills mismatches, geographic mismatch between
vacancies and workers, and pandemic-related labour market restructuring.
By 2024--25, the curve had shifted partially back toward its pre-COVID
position as labour markets adjusted.}
\label{fig:beveridge}
\source{Statistics Canada, Job Vacancy and Wage Survey (Table 14-10-0325-01);
Labour Force Survey (Table 14-10-0287-01). Year labels show last two digits.}
\end{figure}

An outward Beveridge shift has specific policy implications: it means
that reducing the unemployment rate by cutting interest rates (stimulating
demand) will not efficiently eliminate the vacancies. The unfilled positions
require workers with skills that the unemployed do not currently possess,
or are located where unemployed workers do not currently live. The
policy solution is \textit{structural}: retraining, apprenticeship
programs, credential recognition for immigrants, and improved geographic
labour mobility --- not monetary stimulus.

%==========================================================
\section{Wage Determination and the Wage-Productivity Gap}
\label{sec:wages}
%==========================================================

%----------------------------------------------------------
\subsection{Competitive Wage Determination}
%----------------------------------------------------------

In a competitive labour market, the equilibrium wage equals the value of
the marginal product of labour:
\begin{equation}
  W^* = P \times \text{MP}_L
  \label{eq:vmp}
\end{equation}

This means wages grow with: (1) increases in the output price $P$ (firms
can pay more when their revenues rise); and (2) increases in labour
productivity $\text{MP}_L$ (more productive workers generate more revenue
per unit of labour). The fundamental prediction is that real wages and
labour productivity should grow together over the long run.

%----------------------------------------------------------
\subsection{The Wage-Productivity Gap: Canada's Most Important Labour Statistic}
%----------------------------------------------------------

The prediction that real wages and labour productivity should grow together
has not been borne out in Canada since the 1970s. Figure~\ref{fig:wage_productivity}
shows that while labour productivity has grown substantially since 1976,
real wages have grown significantly less.

\begin{figure}[H]
\centering
\begin{tikzpicture}
\begin{axis}[
  width=14cm, height=7.5cm,
  xlabel={Year},
  ylabel={Index (1976 = 100)},
  xlabel style={font=\small},
  ylabel style={font=\small},
  xmin=1975, xmax=2026,
  ymin=90, ymax=195,
  xtick={1980,1985,1990,1995,2000,2005,2010,2015,2020,2025},
  xticklabel style={font=\small},
  ytick={100,120,140,160,180,200},
  yticklabel style={font=\small},
  ymajorgrids=true,
  grid style={dashed, gray!20},
  legend pos=north west,
  legend style={font=\small},
  clip=false
]
% Labour productivity index
\addplot[policyblue, thick, mark=none] coordinates {
  (1976,100)(1980,105)(1985,112)(1990,119)(1995,128)
  (2000,142)(2005,152)(2010,157)(2015,163)(2020,168)(2022,170)(2024,173)(2025,175)
};
\addlegendentry{Labour productivity (real output per hour worked)}
% Real wages index
\addplot[casered, thick, dashed, mark=none] coordinates {
  (1976,100)(1980,101)(1985,99)(1990,104)(1995,102)
  (2000,109)(2005,116)(2010,120)(2015,125)(2020,132)(2022,129)(2024,131)(2025,132)
};
\addlegendentry{Real average weekly earnings}
% Gap annotation
\draw[<->, gray!60, thick] (axis cs:2024,131) -- (axis cs:2024,173);
\node[right, font=\footnotesize, color=gray!60, align=left] at (axis cs:2024.5,152)
  {Gap:\\$\approx$42 pts};
\node[font=\footnotesize, policyblue] at (axis cs:2000,147)
  {Productivity};
\node[font=\footnotesize, casered] at (axis cs:2000,107)
  {Real wages};
\end{axis}
\end{tikzpicture}
\caption{Real Wages vs.\ Labour Productivity, Canada, 1976--2025 (Index 1976=100).
Labour productivity --- real output per hour worked --- has grown by approximately
75\% since 1976. Real average weekly earnings have grown by only approximately
32\% over the same period. The divergence began in the early 1980s and has
persisted, implying that workers have captured a shrinking share of the gains
from productivity growth. This gap is one of the most important distributional
facts in the Canadian economy. It is associated with declining union coverage,
rising profit margins, the growth of monopsony power, and the shift from goods
to services employment --- all discussed below.}
\label{fig:wage_productivity}
\source{Statistics Canada, Table 36-10-0222-01 (multifactor productivity);
Table 14-10-0064-01 (average weekly earnings); author index calculations.
Real wages deflated using CPI (2002=100).}
\end{figure}

The wage-productivity gap implies that the \term{labour share of income}
--- wages and salaries as a fraction of national income --- has declined
over the past four decades, while the \term{capital share} (profits,
dividends, interest) has risen. This matters for inequality (explored
in Chapter~\ref{ch:inequality}) because capital income is far more
concentrated among high-income households than labour income.

The causes of the wage-productivity gap are contested among economists,
but the leading explanations include:

\begin{enumerate}
  \item \textbf{Declining unionization:} Unions compress wage distributions
        and raise wages for medium-skill workers through collective bargaining.
        The decline in union coverage described in Section~\ref{sec:unions}
        has removed a key mechanism for translating productivity gains into
        wage gains.
  \item \textbf{Rising market power of employers (monopsony):} As large
        firms and platforms grow, their ability to set wages unilaterally
        rather than take market wages increases --- widening the gap between
        $W$ and $\text{VMP}_L$.
  \item \textbf{Trade and offshoring:} Competitive pressure from lower-wage
        countries constrains wage growth in tradable sectors, particularly
        manufacturing.
  \item \textbf{Technological change:} Automation has raised the marginal
        product of capital relative to labour in some sectors, shifting
        factor income shares toward capital.
  \item \textbf{Composition effects:} The shift from high-wage manufacturing
        to lower-wage services employment changes the average measured wage
        even if wages within each sector are growing normally.
\end{enumerate}

%==========================================================
\section{Unionization and Collective Bargaining}
\label{sec:unions}
%==========================================================

A \term{trade union} is a collective organisation of workers that bargains
with employers on behalf of its members over wages, benefits, hours, and
working conditions. \term{Collective bargaining} replaces individual
wage negotiation with a negotiation between the union (representing all
workers in a bargaining unit) and the employer. The union's market power
derives from the credible threat of collective action --- a strike.

\begin{figure}[H]
\centering
\begin{tikzpicture}
\begin{axis}[
  ybar,
  bar width=0.35cm,
  width=13cm, height=7cm,
  symbolic x coords={{All industries},{Public admin},{Education},{Healthcare},
                     {Utilities},{Construction},{Manufacturing},
                     {Finance},{Retail},{Food service}},
  xtick=data,
  x tick label style={rotate=40, anchor=east, font=\small},
  ylabel={Union Coverage Rate (\%)},
  ylabel style={font=\small},
  ymin=0, ymax=85,
  ytick={0,10,20,30,40,50,60,70,80},
  yticklabel style={font=\small},
  ymajorgrids=true,
  grid style={dashed, gray!25},
  enlarge x limits=0.05,
  legend pos=north east,
  legend style={font=\small},
  clip=false
]
% 2004 bars
\addplot[fill=policyblue!70, draw=policyblue!90] coordinates {
  ({All industries},32.2)({Public admin},74.3)({Education},72.0)({Healthcare},57.2)
  ({Utilities},59.8)({Construction},31.4)({Manufacturing},30.5)
  ({Finance},8.5)({Retail},13.8)({Food service},7.2)
};
\addlegendentry{2004}
% 2024 bars
\addplot[fill=casered!60, draw=casered!80] coordinates {
  ({All industries},27.8)({Public admin},73.1)({Education},70.5)({Healthcare},54.3)
  ({Utilities},57.2)({Construction},29.8)({Manufacturing},22.4)
  ({Finance},6.1)({Retail},10.2)({Food service},5.8)
};
\addlegendentry{2024}
\end{axis}
\end{tikzpicture}
\caption{Union Coverage Rate by Industry, Canada, 2004 vs.\ 2024.
Unions are heavily concentrated in the public sector --- government,
education, healthcare, and utilities --- where coverage exceeds 54--73\%.
Private-sector union coverage is far lower and declining: manufacturing
fell from 30.5\% to 22.4\%; retail from 13.8\% to 10.2\%. The most
dramatic decline has been in the private sector overall, where coverage
has fallen to approximately 14\%. This bifurcated structure --- heavily
unionized public sector, weakly unionized private sector --- has important
distributional consequences: public-sector workers have maintained real
wages and benefit coverage while many private-sector workers (especially
in retail and food service) have seen stagnant real wages over the same period.}
\label{fig:unionization}
\source{Statistics Canada, Labour Force Survey, Union Status supplement;
Workplace and Employee Survey. 2024 values estimated from recent LFS data.}
\end{figure}

The economic evidence on unions' effects in Canada is reasonably clear:

\begin{itemize}
  \item \textbf{Union wage premium:} Unionized workers earn approximately
        15--20\% more than comparable non-unionized workers in the same
        industry and occupation, after controlling for observable differences.
        The premium is larger for workers without post-secondary education.
  \item \textbf{Benefit premium:} Unionized workers are significantly more
        likely to have employer-sponsored pension plans, drug benefits, and
        dental coverage.
  \item \textbf{Wage compression:} Unions tend to compress the wage
        distribution, raising wages at the bottom more than at the top.
        As union coverage has declined, within-industry wage inequality
        has increased.
  \item \textbf{Productivity effects:} The evidence on unions and productivity
        is mixed. The ``collective voice'' hypothesis suggests unions improve
        productivity by reducing turnover and providing a mechanism for workers
        to communicate grievances. The ``monopoly'' hypothesis suggests unions
        create inefficiencies by raising wages above competitive levels.
        On balance, Canadian studies find small positive or neutral productivity
        effects for most unionized firms.
\end{itemize}

%==========================================================
\section{Minimum Wages in Canada}
\label{sec:min_wage}
%==========================================================

\begin{defbox}{The Minimum Wage}
A \textbf{minimum wage} is a legislated floor on the hourly wage that
employers may legally pay. In Canada, minimum wages are set provincially,
creating significant variation across jurisdictions. There is also a
federal minimum wage covering workers in federally regulated industries
(banking, telecommunications, air transport, interprovincial trucking).

\textit{Economic theory predicts opposite effects under different market
structures:}
\begin{itemize}
  \item \textbf{Competitive market:} A minimum wage above the competitive
        equilibrium causes unemployment (excess labour supply). This is
        Panel A of Figure~\ref{fig:labour_markets}.
  \item \textbf{Monopsony:} A minimum wage up to the competitive wage
        level increases both wages and employment. Panel B of
        Figure~\ref{fig:labour_markets}.
\end{itemize}

The \textit{empirical} evidence (summarised below) is closest to the
monopsony prediction: moderate minimum wage increases in Canada and the
US have generally had small or negligible employment effects, with
significant wage gains for affected workers.
\end{defbox}

\begin{figure}[H]
\centering
\begin{tikzpicture}
\begin{axis}[
  ybar,
  bar width=0.55cm,
  width=13cm, height=6.5cm,
  symbolic x coords={{NL},{NS},{NB},{PEI},{QC},{ON},{MB},{SK},{AB},{BC},{Federal}},
  xtick=data,
  x tick label style={font=\small},
  ylabel={Minimum Wage (\$ per hour), 2025},
  ylabel style={font=\small},
  ymin=13, ymax=19,
  ytick={13,14,15,16,17,18,19},
  yticklabel style={font=\small},
  ymajorgrids=true,
  grid style={dashed, gray!25},
  enlarge x limits=0.06,
  nodes near coords,
  nodes near coords style={font=\footnotesize},
  every axis plot/.append style={fill=unbcgreen!65, draw=unbcgreen!85}
]
\addplot coordinates {
  ({NL},  15.60)({NS},  15.70)({NB},  15.30)({PEI}, 16.00)({QC},  16.10)
  ({ON},  17.20)({MB},  15.80)({SK},  15.00)({AB},  15.00)({BC},  17.40)({Federal},17.30)
};
\end{axis}
\end{tikzpicture}
\caption{Provincial and Federal Minimum Wages, Canada, 2025.
British Columbia (\$17.40) and Ontario (\$17.20) have the highest minimum
wages; Alberta and Saskatchewan (\$15.00) the lowest. The federal minimum
wage (\$17.30) applies to workers in federally regulated industries and
since 2021 is indexed to the Consumer Price Index, preventing real erosion.
The range of \$15.00--\$17.40 represents a 16\% spread across provincial
jurisdictions --- significant enough to affect labour market decisions near
provincial borders and in industries where workers are geographically mobile.}
\label{fig:min_wages}
\source{Employment and Social Development Canada; provincial governments.
Rates as of mid-2025. Some provinces index automatically to CPI.}
\end{figure}

\begin{canexample}{Card and Krueger (1994) and the Canadian Evidence}
David Card and Alan Krueger's 1994 study compared fast-food employment
in New Jersey (which raised its minimum wage from \$4.25 to \$5.05) to
neighbouring Pennsylvania (no change). They found \textit{no significant
employment loss} in New Jersey --- contradicting the simple competitive
model's prediction and consistent with the monopsony model.

Card received the 2021 Nobel Prize in Economics, partly for this work.

Canadian evidence broadly supports the Card-Krueger finding for moderate
increases. A 2019 study by Dube, Lindner, and Reich using Canadian
province-level data found that a 10\% minimum wage increase reduced
employment in the lowest-wage sectors by approximately 1--3\% ---
a small effect relative to the wage gain. A 2021 PBO analysis found
that Ontario's 2018 increase to \$14/hour (a 21\% jump) had modest
negative employment effects of roughly 0.2--0.4\% --- again, small
relative to the wage increase achieved.

However, these averages conceal important heterogeneity. Small businesses
in rural, low-income markets face less monopsony power and therefore
exhibit larger negative employment effects from minimum wage increases
than large firms in urban markets. The ``right'' minimum wage is not
uniform across the vast diversity of Canadian labour markets.
\end{canexample}

%==========================================================
\section{Non-Standard Employment and the Gig Economy}
\label{sec:nonstandard}
%==========================================================

The standard employment relationship --- full-time, permanent, with a single
employer, receiving benefits and entitled to all employment standards --- has
been gradually replaced in parts of the Canadian economy by an increasingly
diverse set of \term{non-standard employment} arrangements.

\begin{figure}[H]
\centering
\begin{tikzpicture}
\begin{axis}[
  width=14cm, height=6.5cm,
  xlabel={Year},
  ylabel={Share of Total Employment (\%)},
  xlabel style={font=\small},
  ylabel style={font=\small},
  xmin=1999, xmax=2025.5,
  ymin=0, ymax=20,
  xtick={2000,2004,2008,2012,2016,2020,2024},
  xticklabel style={font=\small},
  ytick={0,5,10,15,20},
  yticklabel style={font=\small},
  ymajorgrids=true,
  grid style={dashed, gray!20},
  legend pos=north east,
  legend style={font=\small},
  clip=false
]
% Own-account self-employment
\addplot[policyblue, thick, mark=none] coordinates {
  (2000,10.5)(2004,10.8)(2008,10.2)(2012,10.5)(2016,10.9)
  (2020,11.4)(2021,10.6)(2022,10.4)(2023,10.6)(2024,10.8)
};
\addlegendentry{Own-account self-employment}
% Temporary employment
\addplot[casered, thick, dashed, mark=none] coordinates {
  (2000,11.8)(2004,12.4)(2008,12.6)(2012,12.7)(2016,13.2)
  (2020,10.4)(2021,12.5)(2022,13.1)(2023,13.5)(2024,13.7)
};
\addlegendentry{Temporary employment}
% Involuntary part-time
\addplot[dataorange, thick, dotted, mark=none] coordinates {
  (2000,4.2)(2004,4.5)(2008,4.0)(2012,5.2)(2016,5.8)
  (2020,6.5)(2021,5.2)(2022,4.2)(2023,4.5)(2024,4.8)
};
\addlegendentry{Involuntary part-time}
\end{axis}
\end{tikzpicture}
\caption{Non-Standard Employment as a Share of Total Employment, Canada,
2000--2024. Three forms of non-standard work are shown: own-account self-employment
(solo workers without employees, including many gig workers); temporary employment
(term contracts, casual, seasonal); and involuntary part-time work (workers who
want full-time hours but can only find part-time). Together, these categories
account for approximately 29\% of total Canadian employment. Temporary employment
shows a cyclical pattern, falling sharply in the COVID recession as these workers
lost jobs first and recovering quickly. Own-account self-employment has grown
modestly as platform work (delivery, ride-share, freelancing) expands.}
\label{fig:nonstandard}
\source{Statistics Canada, Labour Force Survey, Table 14-10-0190-01. Annual averages.}
\end{figure}

\begin{table}[H]
\centering
\caption{Standard vs.\ Non-Standard Employment: Wages, Benefits, and Security}
\label{tab:employment_comparison}
\renewcommand{\arraystretch}{1.30}
\begin{tabular}{lrrrr}
\toprule
\textbf{Measure} & \textbf{Permanent} & \textbf{Temporary} & \textbf{Own-Account} & \textbf{Platform} \\
                 & \textbf{Full-Time} & \textbf{Employment} & \textbf{Self-Empl.} & \textbf{(Gig)} \\
\midrule
Median hourly wage (\$)   & 31.50 & 22.40 & 24.20 & 18.50 \\
Employer pension (\%)     &  55.3 &  14.2 &   0.0 &   0.0 \\
Employer health benefits  &  64.1 &  19.8 &   0.0 &   0.0 \\
EI eligible (\%)          &  98.0 &  75.0 &   2.5 &   2.5 \\
Paid vacation (\%)        &  97.5 &  43.2 &   0.0 &   0.0 \\
Job tenure (median yrs)   &   7.2 &   1.3 &   5.8 &   1.1 \\
Union coverage (\%)       &  30.2 &   9.5 &   1.1 &   0.5 \\
\bottomrule
\end{tabular}
\source{Statistics Canada, Survey on Employment and Employee Compensation;
Gig Economy Study, Institute for Work and Health, 2024.
Platform/gig values are estimates; sector varies widely.}
\end{table}

Table~\ref{tab:employment_comparison} illustrates the profound gap between
standard and non-standard employment on every dimension: wages, benefits,
Employment Insurance eligibility, and vacation entitlements. This gap has
grown as the gig economy has expanded, and it has concentrated among
younger workers, recent immigrants, and workers without post-secondary
credentials --- precisely the groups with the least bargaining power.

\begin{thinkbox}
Are gig workers ``entrepreneurs who choose flexibility'' or ``misclassified
employees who lack protection''? The answer matters enormously for policy.
If gig workers are choosing flexibility, reclassifying them as employees
could destroy the business model that makes their income possible. If they
are misclassified employees, the current classification denies them legal
rights they are entitled to under employment standards law.

The economic evidence suggests both are partially true: some platform
workers genuinely prefer flexible arrangements, while others are platform
workers because traditional employment is not available to them. A
nuanced policy response --- as adopted in the UK (where Uber drivers
were reclassified as ``workers,'' a middle category with some but not all
employee protections) --- may serve this diversity better than a binary
classification.
\end{thinkbox}

%==========================================================
\section{Demographic Change and the Labour Market}
\label{sec:demographics}
%==========================================================

Canada's labour market is being reshaped by a demographic transition
that was decades in the making. The \term{baby boom generation} (born
1946--1964) is now retiring at the rate of approximately 450{,}000 persons
per year. The cohorts entering the labour force behind them are smaller
(fertility has been below replacement since the 1970s) and less numerous.
The result is a structural tightening of the labour supply that
complements --- and sometimes competes with --- immigration as the
primary mechanism for maintaining the labour force.

\begin{defbox}{The Dependency Ratio}
The \textbf{old-age dependency ratio} measures the number of persons aged
65 and over per 100 persons of working age (15--64):
\begin{equation}
  \text{OADR} = \frac{\text{Population}_{65+}}{\text{Population}_{15-64}} \times 100
  \label{eq:oadr}
\end{equation}

As the baby boom retires and life expectancy increases, Canada's OADR is
rising rapidly: from approximately 22.5 in 2010, to 28.2 today, and
projected to reach 40--45 by 2050 under current fertility assumptions.
Each retiree withdraws from the labour supply (reducing $LF$) while
simultaneously increasing demand for government services (healthcare,
OAS/GIS) that must be financed by the remaining working-age population.

High immigration partially offsets OADR growth because immigrants are
typically working-age adults, but the offset is incomplete: at current
immigration levels, the OADR will still reach approximately 35--38 by 2050,
compared to 40--45 under zero immigration.
\end{defbox}

The demographic transition creates several distinct labour market effects:

\begin{enumerate}
  \item \textbf{Labour shortages in care sectors:} The aging population
        dramatically increases demand for healthcare workers, personal support
        workers, and long-term care staff --- precisely the occupations facing
        the most severe shortages. Canada needs to hire approximately 1 million
        healthcare workers over the next decade simply to maintain current
        service levels for a growing senior population.

  \item \textbf{Skilled trades shortages:} Construction and infrastructure
        maintenance require skilled tradespeople whose training pipeline
        (apprenticeships, technical schools) has consistently underproduced.
        The housing expansion needed to address the affordability crisis
        (Chapter~\ref{ch:housing}) requires more construction workers than
        Canada is currently training.

  \item \textbf{Loss of experience:} When a baby boom machinist like Jerome
        retires, he takes decades of tacit knowledge with him. The
        ``experience premium'' --- the wage differential that rewards years of
        accumulated skill --- is real and measurable. Replacing experienced
        workers with younger or immigrant workers takes years of on-the-job
        learning to replicate.

  \item \textbf{Upward pressure on wages:} As labour becomes relatively
        scarcer, the bargaining power of workers increases. The 2021--2022
        labour market tightness produced the strongest wage growth in
        decades --- particularly in trades, transport, and care work.
        The aging-driven tightness is likely structural rather than cyclical:
        it will persist even as the Bank of Canada slows the economy.
\end{enumerate}

%==========================================================
\section{Case Studies}
\label{sec:ch4_casestudies}
%==========================================================

\begin{casestudy}{The Canadian Labour Market Through COVID-19 and Beyond}

\textbf{Background:} No event in recent Canadian history tested labour
market institutions more severely than the COVID-19 pandemic. The shutdown
of April 2020 --- when over 2 million jobs disappeared in a single month ---
was followed by one of the fastest recoveries in Canadian economic history.
Understanding why illuminates both the resilience and the fault lines
of Canada's labour market.

\textbf{The Collapse:} In April 2020, the unemployment rate hit 13.7\% ---
the highest since the Great Depression. But this understates the disruption:
millions of Canadians who were still employed worked zero hours (treated
as employed-but-absent by the LFS), and the participation rate fell to
62.4\% as many stopped searching. The \textit{actual} share of Canadians
with zero employment income in April 2020 was closer to 25\%.

The collapse was highly concentrated in specific sectors: accommodation
and food services (lost over 40\% of employment); arts and entertainment
(over 35\%); retail trade (over 20\%). These sectors share a key feature:
they require in-person contact and cannot be substituted with remote work.
They also share a demographic profile: disproportionately young, female,
and low-wage workers. The COVID recession was dramatically unequal.

\textbf{The Recovery:} Canada's employment recovery was faster than most
advanced economies. By September 2021, total employment had fully recovered
its pre-pandemic level. By mid-2022, the unemployment rate had fallen to
4.9\% --- a 50-year low. Three factors explain the speed of recovery:
(1) federal income supports (CERB, CEWS) prevented permanent worker
detachment from employers; (2) pent-up consumer demand was released rapidly
when restrictions lifted; and (3) the Bank of Canada maintained
ultra-accommodative monetary conditions through 2021.

\textbf{The Residual:} Full recovery in aggregate conceals persistent
scars. Participation rates for some groups --- particularly prime-age
women (25--54) who took on caregiving responsibilities during school
closures --- remained below pre-COVID levels into 2023. Long-COVID
has kept an estimated 50{,}000--100{,}000 Canadians out of the
labour force. And the Beveridge Curve's outward shift (Figure~\ref{fig:beveridge})
shows that the labour market is less efficient at matching than before
the pandemic --- a structural consequence of the COVID-era disruption
that will take years to fully unwind.
\end{casestudy}

%----------------------------------------------------------
\begin{casestudy}{The Gig Economy and Platform Work in Canada}

\textbf{Background:} Between 2015 and 2024, the number of Canadians
earning income primarily through digital platforms --- Uber, Lyft, DoorDash,
Instacart, TaskRabbit, and similar apps --- grew from negligible to
approximately 800{,}000 to 1{,}000{,}000 workers, representing roughly
4--5\% of the employed workforce. Platform work has attracted enormous
policy attention because it straddles the line between employment and
self-employment in ways that existing labour law was not designed to handle.

\textbf{The Classification Problem:} Under Canadian provincial employment
standards legislation, the key distinction is between \textit{employees}
(entitled to minimum wage, overtime, paid vacation, EI access, and workplace
protections) and \textit{independent contractors} (responsible for their
own taxes, benefits, and business expenses, but generally free to set
their own terms). Platforms classify their workers as independent contractors,
arguing that the flexibility and algorithmic matching that platforms provide
constitutes a genuine market intermediary relationship rather than employment.

\textbf{Workers' perspective:} Platform workers report an effective hourly
wage that, after deducting vehicle costs, fuel, insurance (higher for
commercial use), phone plans, and other business expenses, frequently falls
below the minimum wage. In Ontario, the minimum wage in 2025 is \$17.20/hour;
studies estimate that delivery workers' net hourly income is typically
\$12--\$16/hour after vehicle expenses.

\textbf{Policy responses:} Ontario's 2021 Bill 88 introduced a ``digital
platform worker'' category with some but not all employee protections:
minimum earnings per hour of work (not per assigned period), transparency
requirements on pay calculation, and the right to dispute deactivations.
It stopped short of full reclassification. California's Proposition 22
(2020) similarly created a third-category status for app-based workers.

\textbf{Economic assessment:} The gig economy creates real value by
improving the utilisation of underused assets (vehicles) and providing
flexible income to people who genuinely prefer episodic work. But it also
creates real externalities: the misclassification of employees as contractors
shifts costs from platforms to workers and to public programs (EI, healthcare,
CPP) funded by other workers and employers. The efficient policy
distinguishes between \textit{genuine} flexible self-employment and
misclassified employment --- a distinction that requires bright-line
legal tests that Canadian employment law has not yet fully developed.
\end{casestudy}

%----------------------------------------------------------
\begin{casestudy}{Northern BC: Labour Shortages at the Margin of the Economy}

\textbf{Background:} The University of Northern British Columbia sits in
a region that illustrates labour market dynamics largely invisible in
national statistics. Northern BC faces simultaneous labour surpluses in
some occupations (laid-off mill workers when lumber markets soften) and
acute shortages in others (healthcare workers, skilled tradespeople,
early childhood educators).

\textbf{The Healthcare Worker Shortage:} Northern Health Authority,
which serves over 300{,}000 people across 600{,}000 square kilometres,
consistently operates with nurse vacancy rates of 15--25\%. The shortage
has several causes: geographic remoteness makes recruitment from southern
Canada or internationally difficult; housing costs in resource communities
have risen dramatically with the LNG construction boom; wage competition
from Lower Mainland health authorities and Alberta offers is intense;
and the combination of demanding work conditions and geographic isolation
produces high turnover.

The economic consequence is both a welfare loss (residents receive
suboptimal healthcare) and a fiscal one (travel medicine and locum
staffing costs are far higher per patient than permanent staff costs).

\textbf{The Trades Shortage:} LNG Canada's construction at Kitimat and
related pipeline infrastructure require more skilled tradespeople than
Northern BC has historically trained. The Northern Opportunities program
--- a partnership between School District 57 (Prince George), UNBC, the
College of New Caledonia, and the Northern Development Initiative Trust
--- attempts to address the shortage by channelling Northern BC youth
into pre-apprenticeship training. The program illustrates a general
principle: effective labour market policy is geographically specific,
sector-specific, and requires coordination across educational, government,
and industry actors.

\textbf{The Policy Lesson:} The Beveridge Curve analysis (Section~\ref{sec:beveridge_curve})
showed that Canada's post-COVID outward shift reflects skills and geographic
mismatches. Northern BC is a microcosm of this: vacancies in healthcare and
trades, potential workers in resource sectors facing cyclical disruption.
The solution is not aggregate demand stimulus but targeted training,
credential recognition, geographic incentives, and housing investment
in underserved communities.
\end{casestudy}

%==========================================================
\begin{policydebate}{Should Canada Raise the Federal Minimum Wage to \$20 Per Hour?}

\textbf{The Issue:} The federal minimum wage (covering workers in federally
regulated industries) reached \$17.30/hour in 2025, indexed to CPI.
Labour advocates propose a target of \$20/hour --- a ``living wage''
benchmark derived from estimates of the income needed to cover basic costs
in major Canadian cities. Business groups oppose the increase, citing
potential job losses and cost pressures on small businesses.

\textbf{The Economic Case For \$20/Hour:}
\begin{itemize}
  \item The federal minimum wage applies to a relatively narrow group
        (banking, telecommunications, airlines, interprovincial transport
        --- fewer than 1 million workers). A \$20 federal minimum would
        have modest aggregate employment effects.
  \item The monopsony evidence (Card, Krueger, and Canadian follow-up studies)
        suggests moderate minimum wage increases have small negative employment
        effects and significant positive wage effects for affected workers.
  \item A \$17.30 minimum is below the estimated cost of living in Vancouver
        (\$22/hour), Toronto (\$21/hour), and other major cities. Real-terms
        adequacy requires periodic above-inflation increases.
  \item Higher minimum wages reduce reliance on government transfers,
        shifting income support costs from taxpayers to employers who
        benefit from low-cost labour.
\end{itemize}

\textbf{The Economic Case Against:}
\begin{itemize}
  \item A 15.6\% increase from \$17.30 to \$20.00 exceeds the range
        supported by the monopsony-compatible evidence base. The studies
        finding small employment effects generally examine 5--15\% increases.
        Larger jumps in competitive markets may produce the job losses
        the standard model predicts.
  \item The federal minimum wage creates competitive disparities between
        federally regulated firms (subject to the higher wage) and provincially
        regulated firms (subject to lower provincial minimums). This can
        distort industry structure.
  \item Small businesses, particularly in lower-wage markets, face genuine
        cash-flow constraints. A rapid increase may force closures rather than
        wage adjustments.
  \item A geographically uniform \$20 minimum applied to rural and
        northern communities with lower cost-of-living may exceed the
        competitive equilibrium wage in those markets, producing real
        employment losses there even if it is appropriate in Toronto.
\end{itemize}

\textbf{The Evidence:}
The Bank of Canada has estimated that each 10\% minimum wage increase
reduces youth employment by approximately 0.5\% and total employment by
approximately 0.2\%, with larger effects in lower-wage markets. The PBO
estimates that the proposed \$20 federal minimum would: raise wages for
approximately 200{,}000--250{,}000 workers; result in net employment effects
of approximately $-$10{,}000 to $-$25{,}000 jobs; and reduce government
expenditure on income support programs by approximately \$300--\$500 million
per year as workers earn above benefit thresholds.

\textbf{The Policy Synthesis:} The economic evidence supports moderate
minimum wage increases to a ``living wage'' benchmark over time, with
CPI indexation to prevent real erosion. A stepped increase from \$17.30
to \$20.00 over three years --- rather than a one-time jump --- would
allow labour markets to adjust more gradually and reduce the risk of
discontinuous employment effects. The geographic heterogeneity of Canadian
labour markets also argues for allowing provincial variation rather than
imposing a single federal standard on communities with very different
wage structures.
\end{policydebate}

%==========================================================
\begin{everydaylife}{Understanding Your Pay Stub: Labour Economics in Practice}

Your pay stub is a document in applied labour economics. Every line
reflects a principle from this chapter.

\textbf{Gross wages} represent your nominal wage times hours worked ---
the payment for your marginal product of labour at the market rate
(competitive) or the rate your employer sets unilaterally (monopsony).

\textbf{CPP deductions} (Canada Pension Plan): You contribute 5.95\%
of eligible earnings, matched by your employer. This is a mandatory
contribution to a defined-benefit public pension --- protecting you
against the risk that you save too little for retirement. Gig workers
classified as self-employed pay both the employee and employer share
(11.9\%).

\textbf{EI deductions} (Employment Insurance): You contribute 1.66\%
of insurable earnings. EI provides income replacement if you become
unemployed through no fault of your own --- a form of insurance against
frictional and cyclical unemployment. Platform workers classified as
contractors do not pay into EI and cannot claim it.

\textbf{Income tax withholding}: The progressive income tax system means
your employer withholds tax at your marginal rate each pay period.
The progressivity of the system (higher earners pay a higher marginal
rate) is a key tool for reducing inequality --- which is why economists
carefully track both pre-tax and after-tax income inequality.

\textbf{Union dues} (if applicable): These fund your union's collective
bargaining, legal services, and strike fund. The union wage premium ---
typically 15--20\% --- is the economic return on this investment.
Non-unionized workers receive no such institutional representation.

What your pay stub does not show: the non-wage benefits your employer
provides (pension contributions, health premiums), the employer's CPP
and EI contributions on your behalf (which equal your own contributions),
and the value of job security, training, and safe working conditions that
vary enormously between standard and non-standard employment relationships.
\end{everydaylife}

%==========================================================
\section*{Chapter Summary}
%==========================================================
\addcontentsline{toc}{section}{Chapter Summary}

\begin{itemize}
  \item In a \textbf{competitive labour market}, the equilibrium wage equals
        the value of the marginal product of labour ($W^* = P \times \text{MP}_L$).
        A minimum wage above $W^*$ creates unemployment. In a
        \textbf{monopsony}, the wage is set below the competitive level and
        employment is restricted; a minimum wage up to $W^*$ can increase
        both wages and employment.

  \item The three key labour force statistics are the
        \textbf{unemployment rate} ($U/LF$), the
        \textbf{labour force participation rate} ($LF/\text{WAP}$), and the
        \textbf{employment rate} ($E/\text{WAP}$). Each captures a different
        dimension of labour market performance; all three are needed for a
        complete picture.

  \item \textbf{Frictional unemployment} (search-related), \textbf{structural
        unemployment} (skills/geographic mismatch), and \textbf{cyclical
        unemployment} (demand deficiency) have different causes and require
        different policy responses. Canada's NAIRU is approximately
        5.5--6.0\%.

  \item The \textbf{Beveridge Curve} plots vacancy rates against unemployment.
        Canada's post-COVID outward shift indicates reduced labour market
        matching efficiency --- a structural problem requiring training and
        credential recognition reforms rather than demand stimulus.

  \item Canada's \textbf{wage-productivity gap} has widened since the 1970s:
        labour productivity has grown ~75\% since 1976 while real wages
        have grown only ~32\%. Declining unionization, rising employer
        market power, and trade competition are leading explanations.

  \item \textbf{Union coverage} in Canada's private sector has declined
        to approximately 14\%, while public-sector coverage remains above
        70\%. The union wage premium of 15--20\% and the benefit premium
        document the real costs of declining private-sector unionization.

  \item \textbf{Non-standard employment} (temporary, own-account
        self-employment, gig/platform) accounts for approximately 29\%
        of Canadian employment and offers significantly lower wages, fewer
        benefits, and less job security than permanent full-time work.

  \item \textbf{Demographic aging} --- baby boom retirements, below-replacement
        fertility, and rising dependency ratios --- is creating structural
        labour shortages in healthcare, skilled trades, and early childhood
        education that persist regardless of the business cycle.

  \item \textbf{Provincial minimum wages} range from \$15.00 (AB, SK) to
        \$17.40 (BC). The empirical evidence from Canada and the US suggests
        that moderate minimum wage increases have small negative employment
        effects and meaningful wage gains for affected workers, particularly
        in monopsonistic labour markets.
\end{itemize}

%==========================================================
\section*{Key Terms}
%==========================================================
\addcontentsline{toc}{section}{Key Terms}

\begin{description}[style=nextline, leftmargin=3.5cm]
  \item[\term{Beveridge Curve}] A plot of the job vacancy rate against
        the unemployment rate; its position reflects labour market
        matching efficiency.
  \item[\term{Collective bargaining}] Negotiation between a union
        representing workers and an employer over wages, benefits, and
        conditions.
  \item[\term{Cyclical unemployment}] Unemployment arising from insufficient
        aggregate demand during an economic downturn.
  \item[\term{Dependency ratio (old-age)}] Population 65+ per 100 persons
        aged 15--64; measures the fiscal and labour-supply burden of aging.
  \item[\term{Discouraged workers}] Persons who want work but have stopped
        actively searching and are therefore excluded from the unemployment
        count.
  \item[\term{Employment rate}] Employed persons as a share of the working-age
        population: $E/\text{WAP} \times 100$.
  \item[\term{Frictional unemployment}] Temporary unemployment arising from
        the normal search-and-matching process.
  \item[\term{Gig economy}] Economic activity mediated by digital platforms
        that match workers with short-term tasks or customers.
  \item[\term{Labour force}] Employed plus unemployed persons; $LF = E + U$.
  \item[\term{Labour Force Participation Rate (LFPR)}] Labour force as a
        share of working-age population: $LF/\text{WAP} \times 100$.
  \item[\term{Labour share of income}] Wages and salaries as a fraction
        of national income; has declined over four decades in Canada.
  \item[\term{Minimum wage}] A legislated floor on hourly wages.
  \item[\term{Monopsony}] A labour market with a single (or dominant)
        employer, who can set wages below the competitive level.
  \item[\term{NAIRU}] Non-Accelerating Inflation Rate of Unemployment;
        the unemployment rate at potential output, consisting of frictional
        and structural components; $\approx 5.5$--$6.0\%$ in Canada.
  \item[\term{Non-standard employment}] Employment arrangements that deviate
        from permanent full-time work: temporary, part-time, self-employed,
        or gig.
  \item[\term{Structural unemployment}] Unemployment arising from skills
        or geographic mismatch between available workers and available jobs.
  \item[\term{Unemployment rate}] Unemployed persons as a share of the
        labour force: $U/LF \times 100$.
  \item[\term{Union wage premium}] The percentage difference in wages between
        comparable unionized and non-unionized workers; approximately 15--20\%
        in Canada.
  \item[\term{Value of the marginal product of labour (VMP$_L$)}] The market
        value of the additional output produced by one more unit of labour:
        $\text{VMP}_L = P \times \text{MP}_L$.
  \item[\term{Wage-productivity gap}] The divergence between real wage growth
        and labour productivity growth since the 1970s.
  \item[\term{Working-age population (WAP)}] All persons aged 15 and over;
        the denominator for the employment rate and LFPR.
\end{description}

%==========================================================
\section*{Review Questions}
%==========================================================
\addcontentsline{toc}{section}{Review Questions}

\begin{enumerate}

\item \textbf{(Concept)} Explain the difference between the competitive
and monopsony models of the labour market. How does each model predict
the employment effect of a minimum wage set above the competitive equilibrium?
Why is the Card-Krueger finding (no employment loss from a moderate minimum
wage increase) consistent with the monopsony model but not the competitive model?

\item \textbf{(Calculation)} In a hypothetical province: working-age
population = 3{,}500{,}000; employed = 2{,}100{,}000; unemployed = 130{,}000;
not in labour force = 1{,}270{,}000.
\begin{enumerate}[(a)]
  \item Calculate the unemployment rate, LFPR, and employment rate.
  \item If 50{,}000 discouraged workers re-enter the labour force and
        20{,}000 find jobs, recalculate all three statistics.
  \item Why did the unemployment rate change in (b) even though 20{,}000
        people found jobs?
\end{enumerate}

\item \textbf{(Types of unemployment)} Classify each of the following
as frictional, structural, or cyclical unemployment, and explain your
reasoning:
\begin{enumerate}[(a)]
  \item A software developer in Waterloo who left her job to find one
        with a higher salary and is currently searching.
  \item A coal mine worker in Nova Scotia who lost his job when the
        mine closed due to the carbon transition.
  \item A retail sales associate in Vancouver who was laid off when
        consumer spending fell sharply during a recession.
  \item A nurse in Prince George who cannot find a job because she
        earned her credentials in the Philippines and they are not
        yet recognised in BC.
\end{enumerate}

\item \textbf{(Beveridge Curve)} Canada's job vacancy rate rose from
2.2\% in 2019 to 4.5\% in 2022 even though the unemployment rate fell
from 5.7\% to 5.3\%. Explain what this means for the Beveridge Curve
and for the interpretation of the 2022 labour market tightness. Is this
evidence of structural or cyclical problems?

\item \textbf{(Wage-productivity gap)} Explain two mechanisms through
which declining union coverage could contribute to a widening
wage-productivity gap. For each mechanism, describe a policy intervention
that could partially reverse it.

\item \textbf{(Policy)} A provincial government is considering requiring
all gig/platform workers to be reclassified as employees. Using the
concepts in this chapter, what would you predict are the effects on:
(a) platform worker wages and benefits; (b) platform company costs and
pricing; (c) the number of gig work opportunities available; and (d) the
measurement of unemployment and LFPR?

\end{enumerate}

%==========================================================
\section*{Quantitative Problems}
%==========================================================
\addcontentsline{toc}{section}{Quantitative Problems}

\begin{enumerate}

\item \textbf{Labour Force Accounting.}
Canada's Labour Force Survey for January 2024 reported: working-age
population 34{,}150{,}000; employed 20{,}580{,}000; unemployed 1{,}380{,}000.
\begin{enumerate}[(a)]
  \item Calculate the labour force, unemployment rate, LFPR, and
        employment rate.
  \item If 200{,}000 discouraged workers re-entered the labour force
        in February 2024 (without finding jobs), recalculate all four
        statistics.
  \item If those same 200{,}000 eventually all found employment by March,
        recalculate again.
  \item Explain the sequence of changes in the unemployment rate across
        these three scenarios. What does this illustrate about the
        limitations of the unemployment rate as a welfare measure?
\end{enumerate}

\item \textbf{Minimum Wage and Employment Effects.}
Suppose a city's fast-food labour market can be described as follows:
\begin{align*}
  L^d &= 15{,}000 - 600W \quad \text{(demand for workers, hours/week)}\\
  L^s &= 3{,}000 + 400W \quad \text{(supply of workers, hours/week)}
\end{align*}
\begin{enumerate}[(a)]
  \item Find the competitive equilibrium wage $W^*$ and employment $L^*$.
  \item If the provincial government sets a minimum wage of \$16.00/hour,
        how many workers are employed? How many are unemployed?
  \item Calculate the unemployment generated by the minimum wage as a
        number and as a percentage of the labour force at the minimum wage.
  \item Now suppose the market is actually monopsonistic. The monopsonist's
        marginal labour cost is $\text{MLC} = -3{,}000 + 5W$... [adapt
        from supply curve: if $W = (L^s - 3{,}000)/400$, then
        $\text{MLC} = (2L^s - 3{,}000)/400$]. Find the monopsony wage
        and employment, and show that the \$16.00 minimum wage in this
        context increases employment rather than decreasing it.
\end{enumerate}

\item \textbf{Wage-Productivity Calculation.}
Use the following Canadian data:

\begin{center}
\begin{tabular}{lrrr}
\toprule
Year & Nominal avg.\ weekly earnings (\$) & CPI (2002=100) & Labour productivity index \\
\midrule
1976 & 230.40 & 38.1 & 100.0 \\
1990 & 543.20 & 78.4 & 119.0 \\
2000 & 685.90 & 95.4 & 142.0 \\
2010 & 889.20 & 116.5 & 157.0 \\
2024 & 1{,}245.80 & 168.2 & 173.0 \\
\bottomrule
\end{tabular}
\end{center}

\begin{enumerate}[(a)]
  \item Calculate real average weekly earnings for each year
        (in 1976 dollars), indexed to 1976 = 100.
  \item Calculate the percentage growth in real wages and in labour
        productivity between 1976 and 2024.
  \item What is the ``wage-productivity gap'' in percentage points?
  \item If real wages had grown at the same rate as labour productivity
        since 1976, what would average real weekly earnings be in 2024
        (in 2024 dollars)?
  \item Using Equation~\ref{eq:vmp}, if labour productivity rose 73\%
        but real wages rose only 32\%, what must have happened to the
        ``$P$'' or to the competitive conditions in the labour market?
        Propose at least two explanations.
\end{enumerate}

\item \textbf{Dependency Ratio Projection.}
Statistics Canada projects the following population for Canada:

\begin{center}
\begin{tabular}{lrrr}
\toprule
Year & Population 65+ (M) & Population 15--64 (M) & OADR \\
\midrule
2020 &  6.6 & 24.0 & ? \\
2030 &  9.3 & 25.5 & ? \\
2040 & 11.1 & 26.0 & ? \\
2050 & 12.2 & 26.8 & ? \\
\bottomrule
\end{tabular}
\end{center}

\begin{enumerate}[(a)]
  \item Calculate the old-age dependency ratio for each year
        using Equation~\ref{eq:oadr}.
  \item If each retiree draws an average of \$25{,}000/year in
        OAS/GIS payments, calculate the total OAS/GIS cost in each year.
  \item If this cost is financed by a payroll tax split equally between
        workers and employers, calculate the implied tax rate (as a share
        of an average annual wage of \$65{,}000) in each year.
  \item What does this imply about the fiscal sustainability of Canada's
        public pension system and the potential need for immigration as
        an offset?
\end{enumerate}

\item \textbf{Non-Standard Employment Calculation.}
Using Table~\ref{tab:employment_comparison}:
\begin{enumerate}[(a)]
  \item Calculate the hourly wage gap (\$) and percentage gap between
        permanent full-time and each of the three non-standard categories.
  \item If a platform/gig worker earns \$18.50/hour but has the following
        deductible costs: vehicle depreciation \$2.40/hour, fuel
        \$1.60/hour, insurance premium \$1.20/hour, phone \$0.40/hour,
        what is the net hourly income? How does this compare to Ontario's
        minimum wage?
  \item A full-time permanent worker earning \$31.50/hour receives employer
        pension contributions of 5\% and health benefits valued at
        \$3.20/hour. What is the total effective compensation?
  \item Calculate the ``total compensation gap'' between the gig worker's
        net income from (b) and the permanent worker's total compensation
        from (c).
\end{enumerate}

\item \textbf{Data Exercise.}
Visit Statistics Canada (\url{www.statcan.gc.ca}) and find the most recent
Labour Force Survey release (Table 14-10-0287-01, seasonally adjusted).
\begin{enumerate}[(a)]
  \item Report Canada's current unemployment rate, employment rate, and LFPR.
  \item Break the unemployment rate down by: age group (15--24, 25--54,
        55+), sex (male, female), and educational attainment.
        Which group has the highest unemployment rate? The lowest?
  \item Find the job vacancy rate from the Job Vacancy and Wage Survey
        (Table 14-10-0325-01). Using Figure~\ref{fig:beveridge}, identify
        where the current data point sits relative to the pre-COVID curve.
        What does this imply about current matching efficiency?
  \item How does Canada's unemployment rate compare to the current US and
        OECD average rates? (Source: OECD.Stat, Labour Force Statistics.)
\end{enumerate}

\end{enumerate}

%==========================================================
\section*{Essay Questions}
%==========================================================
\addcontentsline{toc}{section}{Essay Questions}

\begin{enumerate}

\item \textbf{(600--800 words)} The vignette that opened this chapter
described Jerome (a unionized machinist) and Kayla (a gig logistics worker)
as inhabiting very different labour markets in the same city. Using the
economics of this chapter --- wage determination, non-standard employment,
monopsony, and the minimum wage debate --- explain the institutional and
market forces that produce such different outcomes for workers in the same
economy. What policies would you recommend to reduce the gap? What trade-offs
do your recommendations involve?

\item \textbf{(500--700 words)} The wage-productivity gap has widened in
Canada since the 1970s. Some economists argue this is primarily due to
declining unionization; others argue it reflects technological change and
globalisation. Present the best argument for each explanation, evaluate the
empirical evidence, and then state which explanation you find more convincing
and why. Your answer should connect explicitly to Figure~\ref{fig:wage_productivity}
and Figure~\ref{fig:unionization}.

\item \textbf{(400--600 words)} Canada's labour market faces two simultaneous
structural challenges: an aging population reducing labour supply in many
sectors, and an outward Beveridge Curve shift suggesting reduced matching
efficiency. Are these problems complementary (both reflecting a failure
of skills supply to meet demand) or in tension (one might offset the other)?
Evaluate the role of immigration as a potential solution to both challenges,
and identify at least one condition under which immigration could worsen
the matching efficiency problem rather than improve it.

\end{enumerate}

%==========================================================
\section*{Discussion Questions}
%==========================================================
\addcontentsline{toc}{section}{Discussion Questions}

\begin{enumerate}

\item Should Canada have a single national minimum wage rather than
10 different provincial minimums? What are the economic arguments for
standardisation? What are the arguments for provincial variation? How
does your answer relate to the monopsony model and the geographic
diversity of Canadian labour markets?

\item Jerome's unionized position pays \$34.20/hour, has a defined-benefit
pension, and includes comprehensive benefits. Kayla's gig work pays
effectively \$12--\$16/hour with no benefits or security. Is this
difference entirely the result of market forces reflecting different
marginal products of labour? Or does institutional variation (unions,
employment classifications) explain part of the gap? If the latter,
is the current distribution of labour market ``rents'' economically
efficient? Is it equitable?

\item Canada's old-age dependency ratio is rising rapidly. The three
main policy responses are: (a) increasing immigration of working-age
adults; (b) raising the retirement age; and (c) increasing the CPP
contribution rate to build up the fund. What are the economic trade-offs
of each approach? Which combination would you recommend, and why?

\item A provincial government is considering legislation to require
platforms (Uber, DoorDash, etc.) to classify all workers performing
more than 20 hours of platform work per month as ``employees.'' Predict
the effects of this policy on: (a) gig worker wages and benefits;
(b) ride-share and delivery prices; (c) total platform employment;
and (d) the composition of the gig workforce. Who wins and who loses?

\end{enumerate}
